Recent observations from the James Webb Space Telescope (JWST) have uncovered
a population of compact, red galaxies in the early universe, formally
designated as ``Little Red Dots'' (LRDs). Their spectra imply the presence of
Supermassive Black Holes (SMBHs) with masses that appear disproportionately
large relative to their host galaxies. To independently verify the masses of
these black holes, Tidal Disruption Events (TDEs) have been proposed as an
observational tool. A TDE occurs when a star passes sufficiently close to a
SMBH and is torn apart by tidal forces, producing a luminous flare as the
debris is accreted. In this problem, you will analyze LRD observational data,
deduce their scaling relations, evaluate the breakdown of standard dynamical
models, and calculate parameters for a future cosmic survey mission.

\medskip\noindent\textbf{Part 1 (28 points): Empirical Laws and Physical Significance}

Assume the central SMBH mass $(M_{\mathrm{BH}})$ and the total stellar mass
$(M_*^{\mathrm{tot}})$ of the host galaxy follow a scaling relation:
\[
\log M_{\mathrm{BH}} = \alpha + \beta \log M_*^{\mathrm{tot}}
\]
Observational data for 10 representative LRDs (masses expressed in
$\mathrm{M}_\odot$) are provided in Table D03.1.

\begin{center}
\textbf{Table D03.1.} \textit{LRD data.}

\medskip
\begin{tabular}{ccccc}
\toprule
LRD ID & $\log M_*^{\mathrm{tot}}$ & $\Delta\log M_*^{\mathrm{tot}}$
       & $\log M_{\mathrm{BH}}$ & $\Delta\log M_{\mathrm{BH}}$ \\
\midrule
1  & 8.33 & 0.24 & 6.93 & 0.14 \\
2  & 8.30 & 0.30 & 6.40 & 0.30 \\
3  & 7.99 & 0.37 & 7.19 & 0.12 \\
4  & 8.96 & 0.46 & 8.47 & 0.02 \\
5  & 8.70 & 0.36 & 6.74 & 0.17 \\
6  & 8.16 & 0.43 & 7.26 & 0.11 \\
7  & 9.01 & 0.43 & 7.37 & 0.03 \\
8  & 8.65 & 0.39 & 7.56 & 0.08 \\
9  & 8.29 & 0.35 & 7.52 & 0.08 \\
10 & 8.65 & 0.41 & 7.86 & 0.06 \\
\bottomrule
\end{tabular}
\end{center}

\begin{description}
\item[(a)] (5 points) Defining $x_i = \log M_{*,i}^{\mathrm{tot}}$,
  $y_i = \log M_{\mathrm{BH},i}$, and the simplified effective uncertainty
  $(\sigma_{\mathrm{eff},i})$ for each data point as
  \[
  \sigma_{\mathrm{eff},i} = \sqrt{(\Delta x_i)^2 + (\Delta y_i)^2}
  \]
  Fill the values of the statistical weights
  $w_i = 1/\sigma_{\mathrm{eff},i}^2$ in your answer sheet.
\item[(b)] (12 points) Plot 10 data points $(x_i,\ y_i)$ with both horizontal
  and vertical error bars on graph paper (mark your graph as ``D03.1b'').
\item[(c)] (7 points) Apply the weighted least-squares (WLS) regression method
  using your calculated weights $w_i$ to determine the central values and
  standard errors for the slope $(\beta \pm \Delta\beta)$ and the intercept
  $(\alpha \pm \Delta\alpha)$. Overlay the resulting best-fit line on your
  graph. For your calculation, you may rely on the standard WLS formulas
  provided below:
  \begin{align*}
    S_w &= \sum w_i, \quad S_x = \sum w_i x_i, \quad S_y = \sum w_i y_i \\
    S_{xx} &= \sum w_i x_i^2, \quad S_{xy} = \sum w_i x_i y_i, \quad
      \Delta = S_w S_{xx} - (S_x)^2 \\
    \beta &= (S_w S_{xy} - S_x S_y)/\Delta, \quad
    \alpha = (S_{xx} S_y - S_x S_{xy})/\Delta, \\
    \Delta\alpha &= \sqrt{S_{xx}/\Delta}, \quad
    \Delta\beta = \sqrt{S_w/\Delta}
  \end{align*}
  % sic: the official paper typesets the radicand of Delta-alpha as "S_x x"
  % (missing braces); it is S_xx, as the solution confirms
  (Note: Because this regression incorporates statistical weights $(w_i)$, the
  resulting best-fit line may visually deviate from what you would expect from
  a standard, unweighted linear fit.)
\item[(d)] (4 points) Select the most accurate answer for each question.
  \begin{description}
  \item[(d.1)] (2 point) The observed positive slope $(\beta > 0)$ implies a
    coevolutionary trend, yet LRDs host black holes that are disproportionately
    massive compared to their host galaxies. What does this combination most
    likely suggest about the early universe?
    \begin{description}
    \item[A.] Black holes outpace their host galaxies (stellar component)
      initially, either by forming from heavy seeds or undergoing extreme early
      accretion.
    \item[B.] Galaxies fully assemble their stellar mass long before the
      central black hole begins its rapid growth.
    \item[C.] Star formation and black hole accretion are strictly locked at an
      identical rate throughout cosmic history.
    \item[D.] Dark matter halos suppress black hole growth until the host
      galaxy reaches a critical mass threshold.
    \end{description}
  \item[(d.2)] (2 points) The data show a large scatter in stellar mass
    estimations $(\Delta\log M_*^{\mathrm{tot}})$ for LRDs. What is the primary
    physical reason for this substantial uncertainty at high redshift
    $(z \approx 5)$?
    \begin{description}
    \item[A.] Interstellar dust significantly absorbs the black hole's X-ray
      emissions, altering the mass-to-light ratio.
    \item[B.] Rapid stellar collisions in the dense galactic core create
      transient luminous events that mimic a heavier stellar mass.
    \item[C.] The James Webb Space Telescope's UV sensors suffer from extreme
      thermal noise when observing at $z \approx 5$.
    \item[D.] The luminous accretion disk outshines the stellar continuum,
      making it difficult to isolate the host galaxy's light.
    \end{description}
  \end{description}
\end{description}

\medskip\noindent\textbf{Part 2 (6 points): Standard Dynamical Model}

Assume the spatial distribution of the host galaxy's stellar mass follows a
Dehnen spherical density profile,
\[
\rho\,(r) = \frac{(3-\gamma)}{4\pi}\, M_*^{\mathrm{tot}}\,
  \frac{a}{r^{\gamma}(r+a)^{4-\gamma}}
\]
where $a$ is the scale radius (a fundamental parameter marking the transition
zone where the mass density profile shifts between its inner and outer
logarithmic slopes) and $\gamma$ defines the inner density slope
$0 < \gamma < 3$. The black hole's radius of influence $(r_h)$ is defined as
the radius where the enclosed stellar mass equals the black hole mass:
$M_*\,(r_h) = M_{\mathrm{BH}}$. The stellar component has an isotropic velocity
dispersion $\sigma(r)$ and $M(r)$ represents the total enclosed mass within
radius $r$.

\begin{description}
\item[(a)] (2 points) Consider a typical galaxy, applying the ``isothermal
  core'' assumption $(\sigma = \mathrm{const})$ to the region
  $r_h \le r \ll$ a leads to $\gamma = 2$. Therefore, the enclosed stellar mass
  is given by
  \[
  M_*\,(r) = M_*^{\mathrm{tot}} \left( \frac{r}{r+a} \right)
  \]
  Derive the expression for $r_h$ in terms of the scale radius a and the mass
  ratio $M_{\mathrm{BH}}/M_*^{\mathrm{tot}}$.
\item[(b)] (4 points) Based on the result from a), calculate the ratio $r_h/a$
  for LRD ID 4 in Table D03.1. Why does the standard ``isothermal core''
  assumption collapse for this object? Select the most accurate explanation:
  \begin{description}
  \item[A.] The ratio indicates the stellar mass vastly outweighs the black
    hole, rendering the Jeans equation inapplicable.
  \item[B.] The ratio is incredibly small, meaning the black hole has no
    gravitational influence on the core.
  \item[C.] The black hole's radius of influence extends deeply into the scale
    radius.
  \item[D.] The ratio proves that the scale radius a is shrinking over time due
    to dark matter interactions.
  \end{description}
\end{description}

\medskip\noindent\textbf{Part 3 (16 points): Destruction Physics and Data Barriers}

\begin{description}
\item[(a)] (7 points) Consider a star of mass $m_*$ and radius $R_*$
  approaching the SMBH at the distance $r \gg R_*$. The tidal force is defined
  as the difference between gravitational forces exerted by the SMBH at the
  star's center and at its surface. The tidal radius $r = r_{\mathrm{t}}$ is
  determined by balancing the tidal force with the star's self-gravitational
  force at its surface. Derive $r_{\mathrm{t}}$.

  If $r_{\mathrm{t}}$ is smaller than the Schwarzschild radius
  $(R_{\mathrm{S}} = 2GM_{\mathrm{BH}}/c^2)$, the star is swallowed intact
  rather than disrupted. The black hole mass at which
  $r_{\mathrm{t}} = R_{\mathrm{S}}$ is defined as the Hills mass,
  $M_{\mathrm{H}}$. Formulate $M_{\mathrm{H}}$ in terms of $m_*,\ R_*$ and
  calculate its numerical value (in $\mathrm{M}_\odot$) for a Sun-like star.
\item[(b)] (4 points) Using your results
  $(\alpha \pm \Delta\alpha, \beta \pm \Delta\beta)$ from question c) in
  Part 1, determine the maximum stellar mass
  $(\log M_{*,\mathrm{max}}^{\mathrm{tot}})$ a galaxy can have while still
  potentially hosting an observable TDE of a Sun-like star (requiring
  $M_{\mathrm{BH}} = M_{\mathrm{H}})$. Then, calculate the range (Min to Max)
  of $\log M_{*,\mathrm{max}}^{\mathrm{tot}}$. Based on the extreme values,
  select the most accurate statement about the predictive reliability of this
  model:
  \begin{description}
  \item[A.] The model is highly reliable. The calculated range perfectly
    brackets the typical masses of observed galaxies hosting supermassive
    black holes.
  \item[B.] The model lacks predictive power. The lower bound is about the
    mass of a small star cluster, and the upper bound exceeds the mass of the
    observable universe.
  \item[C.] The model fails entirely because a linear weighted least-squares
    regression cannot be mathematically applied to logarithmic astronomical
    data.
  \item[D.] The model is reliable, but the unphysical lower bound indicates
    that we must manually adjust the scale radius to correct the data scatter.
  \end{description}
\item[(c)] (5 points) Cosmic Mission Design

  Assume an all-sky cosmological population of $N \approx 1.021 \times 10^8$
  LRDs located around redshift $z = 5$. The intrinsic (rest-frame) TDE rate per
  galaxy is estimated to be
  $\langle \dot{N}_{\mathrm{TDE}} \rangle \approx 2.01 \times 10^{-6}\,\mathrm{yr^{-1}}$.
  Assume the number of detected TDEs from LRDs to be $5$ during $1$ year of
  continuous observation, what minimum fraction (in percentage) of the
  celestial sphere must the survey area cover?
\end{description}
